\section{Setting, Threat Model, and Positioning}
\label{sec:setting}

\subsection{The epistemic attack surface}
\label{sec:taxonomy}
We organize attacks on the evidential basis of agent decisions by a
ten-class taxonomy of primary attack or failure mechanisms:
misinformation and disinformation; hallucination and factuality
failure; prompt injection and jailbreak; provenance or identity
tampering; retrieval and knowledge poisoning; data or model
poisoning and backdoors; confidence or trust manipulation; tool or
agent hijacking (including MCP tool-description poisoning); drift,
staleness, or benign reliability failure; and a non-attack class for
defense-only, benchmark-only, and evaluation records. The taxonomy was
produced by a registered PRISMA-ScR scoping review~\cite{prisma} of
published work on this surface (6,512 records across six databases;
552 retained as core evidence) whose automated labeling was audited
against two blinded domain experts and an adjudicator on a
hash-registered 20\,\% sample ($n{=}114$): class agreement $\kappa =
0.63$ / Gwet's $AC_1 = 0.65$ between experts, and $\kappa = 0.48$
between experts and pipeline against the pipeline's self-agreement of
0.83. Protocol, ledger, codebook, instruments, hashes, and amendment
log ship with the artifact~\cite{artifact}. The review contributes exactly two things to
this paper: the classes that the crosswalk of
\S\ref{sec:falsifiability} maps onto invariants, and the corpus
prevalence of each class (Appendix~\ref{app:crosswalk}), which is what lets
us say which invariants matter most often. Nothing else in the paper
depends on it.

Categories are mutually exclusive by fiat: where mechanisms co-occur,
the record's own framing decides. One registered disambiguation matters
for the formalization: \emph{retrieval and knowledge poisoning} requires
that corrupted content \emph{enter} the corpus, index, or memory, which
distinguishes substrate corruption (a provenance failure, \Ione) from
manipulation of which authentic evidence is selected, ranked, or
surfaced from an intact substrate (a sufficiency failure, \Ithree); the
latter is carried as a boundary case in the crosswalk.

\subsection{Threat model}
\label{sec:threat}
The adversary controls the untrusted channels: they may inject
instructions or content into retrieved documents and tool results
(indirect prompt injection), poison tool descriptions or servers (MCP
tool poisoning), write to persistent memory or retrieval stores
(memory and RAG poisoning), and shape which evidence the agent sees
(context-selection manipulation). The adversary does not control the
epistemic TCB of \S\ref{sec:tcbdef} or hold its keys. Consistent with
the zero-trust posture now advocated for internal agent
deployments~\cite{controlroadmap}, the model itself is treated as
potentially compromised: a successful injection makes the model an
adversarial component, and every guarantee below must survive that.

Two attack families deserve explicit placement because they live inside
the model rather than on a channel. \emph{Reasoning manipulation and
chain-of-thought poisoning}, in which injected content steers the
model's intermediate reasoning without appearing in any cited evidence,
are, in our terms, attacks on an untrusted component: DEI makes no claim
about the model's internal process, so such attacks surface either as
violations of the counterfactual bound (\Ifive, when removing the
injected span changes the decision) or not at all (when the decision
would have been the same). \emph{Parameter poisoning and backdoors}
triggered by trusted input are declared out of scope: a backdoor fired
by a legitimate directive yields an invariant-satisfying wrong action.
The gateway's own attack surface---what an adversary can do \emph{to}
the monitor rather than around it---is analyzed in \S\ref{sec:security}.

\subsection{Positioning}
\label{sec:positioning}
Out-of-band defenses in the CaMeL lineage~\cite{camel,oob} and
sink-authorization gates such as CXI~\cite{cxi} constrain what a
proposal may \emph{do}; we constrain what may count as a \emph{reason}
for it and prove which part of that constraint a monitor can enforce.
Intent-to-execution integrity~\cite{qu} states the invariance
condition we adopt as \Ifive{} but stops short of a monitor, a
monitorability analysis, or a falsification protocol; agent epistemic
integrity~\cite{shen} prescribes an architecture. Table~\ref{tab:positioning}
records the comparison; \S\ref{sec:related} gives the landscape.

\begin{table}[t]
\centering
\caption{Positioning against the closest prior work.
$\tick$ = provided; $\partl$ = partial; --- = absent.
\emph{Split} = runtime-enforceable separated from auditable-only
guarantees, with supporting theory. All rows verified against the
public record on 17 July 2026.}
\label{tab:positioning}
\scriptsize
\setlength{\tabcolsep}{3pt}
\begin{tabularx}{\columnwidth}{@{}>{\raggedright\arraybackslash}X c c c c c >{\centering\arraybackslash}p{1.05cm}@{}}
\toprule
& Def. & Falsif. & Monit. & Split & Runtime & Dec.-layer \\
\midrule
CaMeL~\cite{camel} & \absent & \absent & \absent & \absent & \tick & \absent \\
GuardAgent~\cite{guardagent} & \absent & \absent & \absent & \absent & \tick & \partl \\
Intent-to-exec.~\cite{qu} & \partl & \absent & \absent & \absent & \absent & \tick \\
Agent epist.\ int.~\cite{shen} & \partl & \absent & \absent & \absent & \absent & \tick \\
CXI~\cite{cxi} & \tick & \absent & \absent & \partl & \tick & \absent \\
ARGUS~\cite{argus} & \absent & \absent & \absent & \absent & \tick & \partl \\
OOB systematization~\cite{oob} & \absent & \partl & \absent & \absent & \absent & \partl \\
\midrule
\textbf{This work} & \tick & \tick & \tick & \tick & \tick & \tick \\
\bottomrule
\end{tabularx}
\end{table}
